\section*{第 \ref{ch:b1:fractions} 章——有理分式}

\begin{solution}{exo:b1:fractions:1}
$\dfrac{1}{X^2 - 1}$：单\omterm{def:b1:fractions:field}{极点} $\pm 1$；遮盖法：
$\dfrac{1/2}{X-1} - \dfrac{1/2}{X+1}$.

$\dfrac{X}{X^2 - 3X + 2} = \dfrac{X}{(X-1)(X-2)}$：遮盖法给出
$\frac{1}{1-2} = -1$ 于 $1$ 处，$\frac{2}{2-1} = 2$ 于 $2$ 处：
$\dfrac{-1}{X-1} + \dfrac{2}{X-2}$.

$\dfrac{X^2+1}{X(X-1)}$：次数差为 $0$，故有整式
部分：作除法，$X^2 + 1 = (X^2 - X) + (X + 1)$，于是 $F = 1 +
\frac{X+1}{X(X-1)}$。对余式用遮盖法，$\frac{1}{-1} = -1$
于 $0$ 处，$\frac{2}{1} = 2$ 于 $1$ 处：
\[
F = 1 - \frac{1}{X} + \frac{2}{X-1} .
\]
\end{solution}

\begin{solution}{exo:b1:fractions:2}
$\dfrac{1}{X(X^2+1)}$：形式为 $\frac aX + \frac{bX + c}{X^2 + 1}$。
在 $0$ 处用遮盖法得 $a = 1$。由 $xF$ 的极限得 $0 = a + b$，故 $b = -1$。
代入 $X = 1$：$\frac12 = 1 + \frac{c - 1}{2}$，故 $c = 0$：
\[
\frac{1}{X(X^2+1)} = \frac 1X - \frac{X}{X^2+1} .
\]

$\dfrac{X^3}{X^2+X+1}$：作除法，$X^3 = (X^2+X+1)(X - 1) + 1$，故
\[
\frac{X^3}{X^2+X+1} = X - 1 + \frac{1}{X^2 + X + 1} ,
\]
已是实分解形式（二次式的判别式
为负）。
\end{solution}

\begin{solution}{exo:b1:fractions:3}
$\dfrac{1}{X^2(X-1)}$：形式为 $\frac{a}{X^2} + \frac bX +
\frac{c}{X-1}$。在二重\omterm{def:b1:fractions:field}{极点} $0$ 处用遮盖法：
$a = \bigl[\frac{1}{X-1}\bigr]_{0} = -1$。在 $1$ 处用遮盖法得 $c = 1$。
由 $xF$ 的极限得 $0 = b + c$，故 $b = -1$：
\[
\frac{1}{X^2(X-1)} = -\frac{1}{X^2} - \frac1X + \frac{1}{X-1} .
\]

$\dfrac{X+1}{(X-1)^3}$：令 $Y = X - 1$，则分子为 $Y + 2$：
\[
\frac{Y + 2}{Y^3} = \frac{1}{Y^2} + \frac{2}{Y^3}
= \frac{1}{(X-1)^2} + \frac{2}{(X-1)^3} .
\]
\end{solution}

\begin{solution}{exo:b1:fractions:4}
\omterm{def:b1:fractions:field}{极点}为四次\omterm{def:b1:complex:unity}{单位根} $1, \iu, -1, -\iu$，均为
单\omterm{def:b1:fractions:field}{极点}。取 $B' = 4X^3$，由单\omterm{def:b1:fractions:field}{极点}公式，$a$ 处的系数为
$\frac{1}{4a^3} = \frac{a}{4a^4} = \frac a4$（利用 $a^4 = 1$）。于是
在 $\C$ 上：
\[
\frac{1}{X^4 - 1}
= \frac{1/4}{X - 1} - \frac{1/4}{X + 1}
+ \frac{\iu/4}{X - \iu} - \frac{\iu/4}{X + \iu} .
\]
合并共轭对（公分母为 $X^2 + 1$）：
$\frac{\iu}{4}\bigl(\frac{1}{X-\iu} - \frac{1}{X+\iu}\bigr) =
\frac{\iu}{4}\cdot\frac{2\iu}{X^2+1} = \frac{-1/2}{X^2+1}$。在
$\R$:
\[
\frac{1}{X^4 - 1}
= \frac{1/4}{X-1} - \frac{1/4}{X+1} - \frac{1/2}{X^2 + 1} .
\]
\emph{在 $X = 0$ 处检验：}$-1 = -\frac14 - \frac14 - \frac12$。
\end{solution}

\begin{solution}{exo:b1:fractions:5}
$\dfrac{X^2}{(X^2+1)^2} = \dfrac{(X^2 + 1) - 1}{(X^2+1)^2} =
\dfrac{1}{X^2+1} - \dfrac{1}{(X^2+1)^2}$.

于是得一个原函数：
\[
\int \frac{x^2\,\dd x}{(x^2+1)^2}
= \arctan x - \frac12\Bigl(\arctan x + \frac{x}{x^2+1}\Bigr) + C
= \frac12\arctan x - \frac{x}{2(x^2+1)} + C .
\]
\end{solution}

\begin{solution}{exo:b1:fractions:6}
\omterm{def:b1:fractions:field}{极点} $0, -1, \dots, -n$ 均为单\omterm{def:b1:fractions:field}{极点}。在 $-k$ 处用遮盖法：
\[
c_k = \frac{n!}{\prod_{j \neq k} (-k + j)}
= \frac{n!}{\bigl(\prod_{j=0}^{k-1}(j - k)\bigr)
\bigl(\prod_{j=k+1}^{n}(j-k)\bigr)}
= \frac{n!}{(-1)^k k!\,(n-k)!} = (-1)^k \binom nk .
\]
于是
\[
\frac{n!}{X(X+1)\cdots(X+n)}
= \sum_{k=0}^{n} \frac{(-1)^k \binom nk}{X + k} .
\]
（检验 $n = 1$ 的情形：$\frac{1}{X(X+1)} = \frac1X -
\frac{1}{X+1}$.)
\end{solution}

\begin{solution}{exo:b1:fractions:7}
$P = X^n - 1$ 有 $n$ 个单根 $\omega^k$
（见\cref{thm:b1:complex:roots}），由
\cref{exo:b1:poly:11}中的对数导数恒等式可得
\[
\sum_{k=0}^{n-1} \frac{1}{X - \omega^k}
= \frac{P'(X)}{P(X)} = \frac{n X^{n-1}}{X^n - 1} .
\]
当 $X = 2$，$n = 4$ 时，右边 $= \frac{4 \times 8}{15} =
\frac{32}{15}$。左边：$\frac{1}{2-1} + \frac{1}{2+1} +
\frac{1}{2 - \iu} + \frac{1}{2 + \iu} = 1 + \frac13 + \frac{4}{5} =
\frac{15 + 5 + 12}{15} = \frac{32}{15}$，利用 $\frac{1}{2-\iu} +
\frac{1}{2+\iu} = \frac{4}{5}$.
\end{solution}

\begin{solution}{exo:b1:fractions:8}
遮盖法：$\dfrac{1}{k(k+1)(k+2)} = \dfrac{1/2}{k} - \dfrac{1}{k+1} +
\dfrac{1/2}{k+2}$。改写为裂项相消的差：
\[
\frac{1}{k(k+1)(k+2)}
= \frac12\Bigl(\frac{1}{k(k+1)} - \frac{1}{(k+1)(k+2)}\Bigr),
\]
（展开验证——或将两个分解式相减）。求和：
\[
S_n = \frac12\Bigl(\frac{1}{1 \times 2} - \frac{1}{(n+1)(n+2)}\Bigr)
= \frac14 - \frac{1}{2(n+1)(n+2)}
\xrightarrow[n \to \infty]{} \frac14 .
\]
\end{solution}

\begin{solution}{exo:b1:fractions:9}
对 $0 \leq m \leq n - 1$，分解分式 $\frac{X^m}{P}$
（次数 $m - n \leq -1$，\omterm{def:b1:fractions:field}{极点}皆为单\omterm{def:b1:fractions:field}{极点}）：由单\omterm{def:b1:fractions:field}{极点}公式得
\[
\frac{X^m}{P} = \sum_{i=1}^{n} \frac{x_i^m / P'(x_i)}{X - x_i} .
\]
乘以 $X$ 并令 $X \to +\infty$：左边趋于
$X^{m+1}/P$ 的极限，其值为 $0$（当 $m \leq n - 2$ 时），为 $1$（当 $m =
n - 1$ 时）（$P$ 为 $n$ 次\omterm{def:b1:poly:def}{首一多项式}）；右边趋于
$\sum_i \frac{x_i^m}{P'(x_i)}$。因此
\[
\sum_{i=1}^{n} \frac{x_i^{m}}{P'(x_i)} =
\begin{cases}
0 & \text{对于 } 0 \leq m \leq n-2,\\
1 & \text{对于 } m = n - 1,
\end{cases}
\]
上式包含了题中所述的两个恒等式（$m = 0$ 要求 $n \geq
2$）。（借助\cref{thm:b1:poly:lagrange}可作如下解释：这些和式是
$X^m$ 的\omterm{thm:b1:poly:lagrange}{拉格朗日插值}\omterm{def:b1:poly:def}{多项式}的首项系数，而
次数 $\leq n-1$ 的\omterm{def:b1:poly:def}{多项式}在 $n$ 个点上插值，
恰好还原该\omterm{def:b1:poly:def}{多项式}。）
\end{solution}

\begin{solution}{exo:b1:fractions:10}
形如 $\frac aX + \frac b{X+1} + \frac c{(X+1)^2}$。遮盖法：在
$0$ 处 $a = 1$；在二重\omterm{def:b1:fractions:field}{极点}处：$c = \bigl[\frac1X
\bigr]_{X=-1} = -1$；由 $xF(x)$ 在无穷远处的极限得 $0 = a + b$，故
$b = -1$：
\[
\frac{1}{X(X+1)^2} = \frac1X - \frac1{X+1} - \frac1{(X+1)^2} .
\]
对 $n = 1, \dots, N$ 求和：前两项裂项相消，得到
$1 - \frac1{N+1}$，第三项贡献 $-\sum_{k=2}^{N+1}
\frac1{k^2}$。令 $N \to \infty$，并利用 $\sum_{k\geq1}
\frac1{k^2} = \frac{\pi^2}6$：
\[
\sum_{n=1}^{\infty}\frac{1}{n(n+1)^2}
= 1 - \Bigl(\frac{\pi^2}6 - 1\Bigr) = 2 - \frac{\pi^2}6
\approx 0.355 .
\]
\end{solution}

\begin{solution}{exo:b1:fractions:11}
在变量 $Y$ 中：$\frac1{(Y+1)(Y+4)} = \frac{1/3}{Y+1} -
\frac{1/3}{Y+4}$（分别在 $-1$ 与 $-4$ 处遮盖），故
\[
\frac{1}{(X^2+1)(X^2+4)}
= \frac13\,\frac1{X^2+1} - \frac13\,\frac1{X^2+4} .
\]
同理 $\frac{Y}{(Y+1)(Y+4)} = \frac{-1/3}{Y+1} +
\frac{4/3}{Y+4}$，故
\[
\frac{X^2}{(X^2+1)(X^2+4)}
= -\frac13\,\frac1{X^2+1} + \frac43\,\frac1{X^2+4} .
\]
（在 $X = 0$ 处检验：$0 = \frac13(-1 + 1)$。）这些已经是
实分解：不可约二次式上的分子
恰好都是常数。
\end{solution}

\begin{solution}{exo:b1:fractions:12}
\begin{enumerate}
  \item 在避开\omterm{def:b1:fractions:field}{极点}的每个区间上，$F'(x) = -\sum_i
        \frac{c_i}{(x - p_i)^2} < 0$，于是严格递减。当 $x
        \to \pm\infty$ 时各项都趋于 $0$，故 $F \to 0$，
        在 $+\infty$ 处从上方趋近（此时各项皆为正），在
        $-\infty$ 处从下方趋近。当 $x \to p_i^+$ 时，
        $\frac{c_i}{x - p_i}$ 趋向 $+\infty$，其余各项
        保持有界，故 $F \to +\infty$；类似地，$F \to -\infty$（当
        $x \to p_i^-$ 时）。
  \item 固定 $\lambda > 0$。在 $\intoo{-\infty}{p_1}$ 上：$F$
        从 $0^-$ 递减至 $-\infty$，故 $F < 0 < \lambda$，
        无解。在每个 $\intoo{p_i}{p_{i+1}}$（$1 \leq i
        \leq r-1$）上：$F$ 从 $+\infty$ 递减至 $-\infty$，
        故恰取值 $\lambda$ 一次（由介值性质及
        严格单调性）。在
        $\intoo{p_r}{+\infty}$ 上：$F$ 从 $+\infty$ 递减至
        $0^+$，同样恰有一个解。总计恰有 $r$
        个解，且与\omterm{def:b1:fractions:field}{极点}相互交错。
\end{enumerate}
\end{solution}

\begin{solution}{pb:b1:fractions:1}
\textbf{1.} 遮盖法：$\frac1{X(X+1)} = \frac1X - \frac1{X+1}$。
求和后裂项相消：
\[
\sum_{n=1}^{N}\Bigl(\frac1n - \frac1{n+1}\Bigr)
= 1 - \frac1{N+1} \longrightarrow 1 .
\]

\textbf{2.} $\frac1{X(X+2)} = \frac{1/2}X - \frac{1/2}{X+2}$.
求和时，$\frac1n$ 项在 $n = 1, 2$ 时保留，而
$-\frac1{n+2}$ 项在 $n = N-1, N$ 时保留：
\[
\sum_{n=1}^{N}\frac1{n(n+2)}
= \frac12\Bigl(1 + \frac12 - \frac1{N+1} - \frac1{N+2}\Bigr)
\longrightarrow \frac34 .
\]

\textbf{3.} 若 $F(X) = G(X) - G(X+1)$，则 $\sum_{n=1}^N F(n) =
\sum_{n=1}^N\bigl(G(n) - G(n+1)\bigr) = G(1) - G(N+1)$：中间各项
两两相消。对于 $F =
\frac1{X(X+1)(X+2)}$，可取 $G(X) = \frac1{2X(X+1)}$：
\[
G(X) - G(X+1)
= \frac{(X+2) - X}{2X(X+1)(X+2)} = F(X) ,
\]
于是 $\sum_{n=1}^N F(n) = \frac14 - \frac1{2(N+1)(N+2)} \to
\frac14$，即\cref{exo:b1:fractions:8}的值。

\textbf{4.} 将右边通分，公分母为
$X(X+1)\cdots(X+k)$：
\[
\frac1k\cdot\frac{(X + k) - X}{X(X+1)\cdots(X+k)}
= \frac{1}{X(X+1)\cdots(X+k)} ,
\]
这正是恒等式。于是 $F_k(X) = \frac1k\bigl(G_k(X) -
G_k(X+1)\bigr)$，其中 $G_k(X) = \frac1{X(X+1)\cdots(X+k-1)}$，由
第 3 题的方法给出
\[
\sum_{n=1}^{N}\frac1{n(n+1)\cdots(n+k)}
= \frac1k\Bigl(\frac1{k!} - G_k(N+1)\Bigr)
\longrightarrow \frac1{k\cdot k!} ,
\]
因为 $G_k(1) = \frac1{k!}$，且 $G_k(N+1) \to 0$。当 $k = 2$ 时：
$\frac1{2\cdot2} = \frac14$，与第 3 题结果一致。

\textbf{5.} 由\cref{exo:b1:fractions:6}可得
$\frac{n!}{X(X+1)\cdots(X+n)} = \sum_{j=0}^n
\frac{(-1)^j\binom nj}{X + j}$。取 $X = 1$：左边
为 $\frac{n!}{(n+1)!} = \frac1{n+1}$，右边为
$\sum_j(-1)^j\binom nj\frac1{1+j}$：第一个恒等式。取 $X = 2$：
左边为 $\frac{n!}{2\cdot3\cdots(n+2)} =
\frac{n!\cdot1}{(n+2)!} = \frac1{(n+1)(n+2)}$，右边为
$\sum_j(-1)^j\binom nj\frac1{2+j}$：第二个恒等式。

\textbf{6.} \omterm{def:b1:fractions:field}{极点} $\omega^k$ 均为单\omterm{def:b1:fractions:field}{极点}，由单\omterm{def:b1:fractions:field}{极点}的
公式（见\cref{met:b1:fractions:compute}）给出系数
\[
\frac{1}{\bigl(nX^{n-1}\bigr)_{X = \omega^k}}
= \frac{1}{n\,\omega^{k(n-1)}}
= \frac{\omega^k}{n\,\omega^{kn}} = \frac{\omega^k}n ,
\]
利用 $\omega^{kn} = 1$。因此 $\frac1{X^n-1} = \frac1n\sum_k
\frac{\omega^k}{X - \omega^k}$。

\textbf{7.} 当 $n = 2$（$\omega = -1$）：$\frac12\bigl(
\frac1{X-1} - \frac1{X+1}\bigr) = \frac12\cdot\frac{2}{X^2-1} =
\frac1{X^2-1}$：正确。各系数之和为
$\frac1n\sum_k\omega^k = 0$，当 $n \geq 2$ 时
（见\cref{prop:b1:complex:sumroots}）。这是必然的：$x\,F(x) \to
\sum_k c_k$，当 $x \to \infty$ 时，对任何只含单
\omterm{def:b1:fractions:field}{极点}的分解都成立，而此处 $xF(x) = \frac{x}{x^n-1} \to 0$，因为 $n \geq
2$。

\textbf{8.} 对 $\frac{X^{n-1}}{X^n - 1}$ 在 $\omega^k$ 处用遮盖法：
$\frac{A(\omega^k)}{B'(\omega^k)} =
\frac{\omega^{k(n-1)}}{n\omega^{k(n-1)}} = \frac1n$，于是
$\frac{X^{n-1}}{X^n-1} = \frac1n\sum_k\frac1{X - \omega^k}$——
这再次是\cref{exo:b1:fractions:7}中的恒等式。

\textbf{9.} 记 $c = \omega^k$，$\conj c = \omega^{n-k}$，并且
$c\conj c = 1$，$c + \conj c = 2\cos\theta_k$：
\[
\frac{c}{X - c} + \frac{\conj c}{X - \conj c}
= \frac{c(X - \conj c) + \conj c(X - c)}
{(X - c)(X - \conj c)}
= \frac{2\cos\theta_k\,X - 2}{X^2 - 2\cos\theta_k\,X + 1} .
\]
在第 6 题中将 $k$ 与 $n - k$ 配对，当 $n$ 为奇数时，
\[
\frac1{X^n - 1} = \frac1n\Biggl(\frac1{X-1} +
\sum_{k=1}^{(n-1)/2}
\frac{2\cos\theta_k X - 2}{X^2 - 2\cos\theta_k X + 1}\Biggr) ;
\]
当 $n$ 为偶数时，多出的自配对\omterm{def:b1:fractions:field}{极点} $\omega^{n/2} = -1$
括号内多出的项为 $\frac{-1}{X+1}$，
配对求和的上限为 $\frac n2 - 1$。

\textbf{10.} $n = 4$：$\theta_1 = \frac\pi2$，$\cos\theta_1 = 0$，
于是配对项为 $\frac{-2}{X^2+1}$，得
\[
\frac1{X^4-1} = \frac14\Bigl(\frac1{X-1} - \frac1{X+1}
- \frac2{X^2+1}\Bigr) ,
\]
即\cref{exo:b1:fractions:4}中的分解式。

\textbf{11.} 令 $P = \prod_{k=1}^{n-1}(X - \omega^k)$（即 $X^n -
1$ 除以 $X - 1$），由\cref{exo:b1:poly:11}，$\frac{P'}P =
\sum_{k=1}^{n-1}\frac1{X-\omega^k}$。取 $X = 1$：$P(1) = n$，且
$P'(1) = \sum_{j=1}^{n-1}j = \frac{n(n-1)}2$，从而
\[
\sum_{k=1}^{n-1}\frac1{1 - \omega^k}
= \frac{P'(1)}{P(1)} = \frac{n-1}2 .
\]

\textbf{12.} 半角公式：$1 - \eu^{\iu\theta} =
-2\iu\sin\frac\theta2\,\eu^{\iu\theta/2}$，于是
\[
\frac1{1 - \eu^{\iu\theta}}
= \frac{\eu^{-\iu\theta/2}}{-2\iu\sin\frac\theta2}
= \frac{\iu\cos\frac\theta2 + \sin\frac\theta2}
{2\sin\frac\theta2}
= \frac12 + \frac\iu2\,
\frac{\cos\frac\theta2}{\sin\frac\theta2} .
\]
令 $\theta = \theta_k = \frac{2k\pi}n$，则 $\frac1{1-\omega^k} =
\frac12 + \frac\iu2\cot\frac{k\pi}n$。对 $k = 1, \dots,
n-1$ 求和，并与第 11 题中的实数值 $\frac{n-1}2$ 比较，
实部已完全吻合，因此
$\sum_k\cot\frac{k\pi}n = 0$——亦可由对称关系 $\cot\frac{(n-k)
\pi}n = -\cot\frac{k\pi}n$ 也可看出。

\textbf{13.} 对 $\frac{P'}P = \sum_k\frac1{X -
\omega^k}$ 求导：
\[
\frac{P''}P - \Bigl(\frac{P'}P\Bigr)^2
= -\sum_{k=1}^{n-1}\frac1{(X - \omega^k)^2} .
\]
当 $X = 1$ 时：$P''(1) = \sum_{j=2}^{n-1}j(j-1) = 2\binom n3 =
\frac{n(n-1)(n-2)}3$（由曲棍球棒恒等式或归纳法），于是
\[
\sum_{k=1}^{n-1}\frac1{(1 - \omega^k)^2}
= \Bigl(\frac{n-1}2\Bigr)^2 - \frac{(n-1)(n-2)}3
= \frac{(n-1)\bigl(3(n-1) - 4(n-2)\bigr)}{12}
= \frac{(n-1)(5-n)}{12} .
\]

\textbf{14.} 将第 12 题的公式平方，记 $c_k =
\cot\frac{k\pi}n$：
\[
\frac1{(1-\omega^k)^2}
= \Bigl(\frac12 + \frac\iu2 c_k\Bigr)^2
= \frac14 - \frac{c_k^2}4 + \frac\iu2\,c_k .
\]
求和，并利用 $\sum c_k = 0$（第 12 题）及第 13 题的结果：
$\frac{n-1}4 - \frac14\sum_k c_k^2 = \frac{(n-1)(5-n)}{12}$，于是
\[
\sum_{k=1}^{n-1}\cot^2\frac{k\pi}n
= (n-1) - \frac{(n-1)(5-n)}3 = \frac{(n-1)(n-2)}3 .
\]
再由 $\frac1{\sin^2t} = 1 + \cot^2t$ 得
$\sum_k\frac1{\sin^2\frac{k\pi}n} = (n-1) +
\frac{(n-1)(n-2)}3 = \frac{n^2-1}3$。

\textbf{15.} 当 $n = 3$ 时：$\cot^2\frac\pi3 + \cot^2\frac{2\pi}3 =
\frac13 + \frac13 = \frac23 = \frac{2\cdot1}3$；又
$\frac1{\sin^2}$ 之和为 $\frac43 + \frac43 = \frac83 =
\frac{9-1}3$。当 $n = 4$ 时：$1 + 0 + 1 = 2 = \frac{3\cdot2}3$；又 $2
+ 1 + 2 = 5 = \frac{16-1}3$。两式均验证无误。

\textbf{16.} 对 $t \in \intoo0{\frac\pi2}$，由经典的
比较式 $\sin t < t < \tan t$（可用面积或凸性论证，
高中分册中已熟悉），取倒数得 $\cot t <
\frac1t < \frac1{\sin t}$，三者在该区间内均为正；两边平方后
不等号方向不变。令 $t = \frac{k\pi}n$，$1 \leq k \leq m$，
$n = 2m+1$（于是 $t < \frac\pi2$）：
\[
\cot^2\frac{k\pi}n < \frac{n^2}{k^2\pi^2} <
\frac1{\sin^2\frac{k\pi}n} .
\]

\textbf{17.} 由对称关系 $\cot^2\frac{(n-k)\pi}n =
\cot^2\frac{k\pi}n$，对 $\sin^2$ 同理，第 14 题的两个和
式减半：$\sum_{k=1}^{m}\cot^2\frac{k\pi}n = \frac{(n-1)(n-2)}6 =
\frac{m(2m-1)}3$，及 $\sum_{k=1}^m\frac1{\sin^2\frac{k\pi}n} =
\frac{n^2-1}6 = \frac{2m(m+1)}3$。对第 16 题在 $k =
1, \dots, m$ 上求和，并乘以 $\frac{\pi^2}{n^2}$：
\[
\frac{\pi^2}{(2m+1)^2}\cdot\frac{m(2m-1)}3
\;<\; \sum_{k=1}^{m}\frac1{k^2} \;<\;
\frac{\pi^2}{(2m+1)^2}\cdot\frac{2m(m+1)}3 .
\]
两界均趋于 $\frac{\pi^2}6$（当 $m \to \infty$ 时，比值
$\frac{m(2m-1)}{(2m+1)^2}$ 与 $\frac{2m(m+1)}{(2m+1)^2}$ 均
趋于 $\frac12$），由夹逼原理，单调递增的部分和
收敛，且
\[
\sum_{k=1}^{\infty}\frac1{k^2} = \frac{\pi^2}6 .
\]

\textbf{18.} 将 $\frac1{X^2-1} = \frac12\bigl(\frac1{X-1} -
\frac1{X+1}\bigr)$ 两边平方：
\[
\frac1{(X^2-1)^2} = \frac14\Bigl(\frac1{(X-1)^2} +
\frac1{(X+1)^2}\Bigr) - \frac12\cdot\frac1{(X-1)(X+1)} ,
\]
再将交叉项重新分解为 $\frac1{(X-1)(X+1)} =
\frac12\bigl(\frac1{X-1} - \frac1{X+1}\bigr)$，即得所述
形式。当 $X = 0$ 时：左边为 $1$；右边为 $\frac14(1 + 1) -
\frac14(-1 - 1) = \frac12 + \frac12 = 1$。

\textbf{19.} 对第 18 题关于 $n \geq 2$ 求和。记 $S =
\sum_{k\geq1}\frac1{k^2} = \frac{\pi^2}6$：$\sum_{n\geq2}
\frac1{(n-1)^2} = S$；$\sum_{n\geq2}\frac1{(n+1)^2} = S - 1 -
\frac14$；又间隔为 2 的裂项相消 $\sum_{n\geq2}\bigl(
\frac1{n-1} - \frac1{n+1}\bigr) = 1 + \frac12 = \frac32$。由此得
\[
\sum_{n=2}^{\infty}\frac1{(n^2-1)^2}
= \frac14\Bigl(2S - \frac54\Bigr) - \frac14\cdot\frac32
= \frac S2 - \frac{11}{16}
= \frac{\pi^2}{12} - \frac{11}{16} \approx 0.135 .
\]
数值上：$\frac19 + \frac1{64} + \frac1{225} + \frac1{576} +
\frac1{1225} + \dots \approx 0.1111 + 0.0156 + 0.0044 + 0.0017 +
0.0008 + \dots \approx 0.135$：与上式一致。

\textbf{20.} 对第 8 题的恒等式求导
$\sum_k\frac1{X-\omega^k} = \frac{nX^{n-1}}{X^n-1}$：
\[
\sum_{k=0}^{n-1}\frac1{(X - \omega^k)^2}
= \frac{n^2X^{2n-2} - n(n-1)X^{n-2}(X^n - 1)}{(X^n - 1)^2} .
\]
验证 $n = 2$，$X = 2$ 时：右边 $\frac{4\cdot4 -
2\cdot1\cdot3}{9} = \frac{10}9$；左边 $\frac1{(2-1)^2} +
\frac1{(2+1)^2} = 1 + \frac19 = \frac{10}9$。

\textbf{21.} $\frac1{\binom nk} = \frac{k!\,(n-k)!}{n!} =
\frac{k!}{(n-k+1)(n-k+2)\cdots n}$，分母是 $k$ 个
连续整数的乘积。令 $j = n - k + 1$（于是 $j$ 取遍
$\N^*$，$n$ 从 $k$ 起）：
\[
\sum_{n=k}^{\infty}\frac1{\binom nk}
= k!\sum_{j=1}^{\infty}\frac1{j(j+1)\cdots(j+k-1)}
= \frac{k!}{(k-1)\,(k-1)!} = \frac{k}{k-1} ,
\]
由第 4 题以 $k - 1$ 代替 $k$ 即得（因
$k - 1 \geq 1$ 成立）。

\textbf{22.} 当 $k = 3$ 时，$\frac1{\binom n3} =
\frac6{(n-2)(n-1)n}$，由第 3 题的部分和（作指标平移）可得
\[
\sum_{n=3}^{N}\frac1{\binom n3}
= 6\sum_{j=1}^{N-2}\frac1{j(j+1)(j+2)}
= 6\Bigl(\frac14 - \frac1{2(N-1)N}\Bigr)
= \frac32 - \frac3{(N-1)N}
\longrightarrow \frac32 ,
\]
即第 21 题所得的 $\frac k{k-1} = \frac32$。

\textbf{23.} 当 $n = 6$ 时，共轭对 $k = 1$（$\theta_1 = \frac\pi3$，
$2\cos\theta_1 = 1$）与 $k = 2$（$\theta_2 = \frac{2\pi}3$，
$2\cos\theta_2 = -1$），再加上实\omterm{def:b1:fractions:field}{极点} $\pm1$：
\[
\frac1{X^6-1} = \frac16\Bigl(\frac1{X-1} - \frac1{X+1}
+ \frac{X - 2}{X^2 - X + 1}
+ \frac{-X - 2}{X^2 + X + 1}\Bigr) .
\]
令 $X = 0$：左边为 $-1$；右边为 $\frac16(-1 - 1 - 2 - 2)
= -1$，验证正确。

\textbf{24.} (i) 唯一性保证了每一次系数比较都有依据
——遮盖法、第 9 题的共轭对分组，以及
求导技巧（第 13、20 题）都依赖于它。
(ii) \omterm{def:b1:complex:unity}{单位根}给出了 $X^n -
1$ 的\omterm{def:b1:fractions:field}{极点}、其对称性（$k \leftrightarrow n-k$），以及
第 12 题中的半角代数，
（见\cref{met:b1:complex:trig}）。(iii) 对数导数
$\frac{P'}P = \sum\frac{m_i}{X-a_i}$ 把关于
$P = 1 + X + \dots + X^{n-1}$ 的根的信息，转化为
第 11、13 题的数值和——第二部分与
第三部分之间的枢纽。

\textbf{25.} 这一分解是一个纯代数恒等式，
对变量的每个取值同时成立；正是这一点使它成为
一台“引擎”。代入整数并求和，把它化为
裂项相消之和（第一部分）；代入\omterm{def:b1:complex:unity}{单位根}并对共轭根分组，把它化为
三角恒等式（第 II、
III 部分）；只有在最后一步，分析才登场——
两个闭形式之间的夹逼，得出 $\frac{\pi^2}6$——
这一结论是任何有限次代入都无法触及的。这种分工
（代数给出精确的有限恒等式，分析负责
取极限）正是 \cref{ch:b1:series} 的模板，其中
裂项相消与比较判别法变得系统化，以及
\cref{ch:b1:integration} 的模板，其中每块“积木”都获得一个原函数，
同样的分解用来计算积分而非求和。
\end{solution}
